\subsection{Groud-truth Dataset Construction}
\label{sec:data-collection}

We evaluate ReCall on web applications with source code publicly available on GitHub. Access to their source code and build configurations allows us to obtain library sources and generate source maps for establishing ground-truth correspondences between source and bundled functions.

We select seven widely used bundlers based on the bundler usage survey\footnote{\label{fn:bundler-usage-survey}State of JS build-tools survey: \url{https://2025.stateofjs.com/en-US/libraries/build-tools/}}: Webpack, Vite, esbuild, Rollup, Turbopack, Parcel, and Bun. For each bundler, we search GitHub for applications that use it and examine repositories in descending order of star count. We require each application to have a frontend page that builds and opens successfully and at least five production library dependencies. This process yields 35 repositories.

For each repository, the ground-truth dataset is consisted of three components: \textbf{library source code}, \textbf{runtime intercepted JavaScript files}, and \textbf{ground-truth function mappings}. Given a repository, we first identify production libraries from project configuration files and randomly select up to ten per repository, yielding 282 library instances in total. Then, we manually enable source-map generation and verify that each application’s frontend opens successfully. Following prior work (CITE LATER), we capture JavaScript files loaded when the home page opens and retrieve their source maps from the build output. Our collection does not cover additional files whose loading requires page interactions.

Before establishing ground-truth mappings, we need to first discuss four special function transformations introduced by the build process that affect recovery: removal, addition, combination, and splitting. Removal eliminates source functions, for example through tree shaking, while addition introduces functions such as generated helpers or functions produced by factories. Removed functions are absent from the bundle, and added functions have no original source counterpart to recover. Recovery is therefore not meaningful in either case, and we exclude both from the recovery evaluation.

Combination and splitting can preserve program behavior while breaking the function-level semantic equivalence on which ReCall relies. To assess their impact on ReCall’s applicability, we estimate the frequency of them using source-map relationship: one-to-many correspondences indicate splitting, many-to-one correspondences indicate combination, and source functions without any mapped bundled counterpart indicate removal. We omit addition because generated library functions are difficult to distinguish from application code.

Our source-map analysis shows that one-to-one correspondences dominate the observed category counts across the 35 repositories. We identify 1,835 split candidates (2.49\%), 1,319 combine candidates (1.79\%), and 12,438 remove candidates (16.87\%), alongside 58,128 one-to-one function pairs (78.85\%). The relatively small shares of split and combine candidates suggest that boundary-changing transformations have a limited impact on ReCall’s applicability. As a result, we can safely construct the ground-truth dataset using \textbf{only one-to-one mappings}, assuming that each bundled function has exactly one correct source-function counterpart in our experiments and vice versa, and leave boundary-changing transformation cases for future work.

Table~\ref{tab:rq1-dataset-bundlers} summarizes the collected dataset which will be used throughout RQ1 to RQ3.

\begin{table}[htbp]
\caption{Ground-truth dataset for library function recovery evaluation, grouped by bundler.}
\label{tab:rq1-dataset-bundlers}
\centering
\small
\renewcommand{\arraystretch}{1.10}
\setlength{\tabcolsep}{3pt}
\begin{tabular*}{0.9\linewidth}{@{\extracolsep{\fill}}lccccc@{}}
\toprule
\textbf{Bundler} &
\multicolumn{1}{c}{\begin{tabular}[c]{@{}c@{}}\textbf{Bundler}\\\textbf{Market Share}\footnotemark\end{tabular}} &
\multicolumn{1}{c}{\begin{tabular}[c]{@{}c@{}}\textbf{\# Collected}\\\textbf{Repos}\end{tabular}} &
\multicolumn{1}{c}{\textbf{\# Lib}} &
\multicolumn{1}{c}{\begin{tabular}[c]{@{}c@{}}\textbf{\# Ground-truth}\\\textbf{Function Pairs}\end{tabular}} &
\multicolumn{1}{c}{\begin{tabular}[c]{@{}c@{}}\textbf{Avg. \# Pairs}\\\textbf{per Lib}\end{tabular}} \\
\midrule
Webpack & 22.89\% & 8 & 75 & 19,055 & 254.1 \\
Vite & 22.11\% & 8 & 79 & 19,630 & 248.5 \\
esbuild & 14.21\% & 6 & 41 & 5,208 & 127.0 \\
Rollup & 12.37\% & 3 & 9 & 1,062 & 118.0 \\
Turbopack & 7.37\% & 6 & 60 & 8,605 & 143.4 \\
Parcel & 6.84\% & 2 & 9 & 2,880 & 320.0 \\
Bun & 6.58\% & 2 & 9 & 1,688 & 187.6 \\
\midrule
\textbf{Total} & \textbf{92.37\%} & \textbf{35} & \textbf{282} & \textbf{58,128} & \textbf{206.1} \\
\bottomrule
\end{tabular*}
\end{table}
\footnotetext{Bundler market shares are approximated by normalizing the usage figures reported in the JS build-tools survey (Footnote~\ref{fn:bundler-usage-survey}). Under this approximation, the seven listed bundlers account for 92.37\% of the market.}



